\chapter{Causing Single-Bit Faults in Practice}
\label{chap:physical_faults}

The attacks in this report assume a very specific event: a short, transient fault
that changes one bit (Giraud) or one byte (Piret) of an intermediate value, at a
chosen instruction, while leaving everything else correct. The mathematics then
turns the correct/faulty ciphertext pair into key information. This chapter
explains, at the level of a laboratory evaluation, how such faults are produced
on real hardware, how far the clean single-bit model is from typical practice,
and which countermeasures remove the attacker's advantage. It is written for a
defensive security evaluation on the project's own ChipWhisperer target; it is
not a procedure against any third-party device.

\section{Why a physical fault appears}
Digital logic is correct only when every signal settles within one clock period
and every transistor is powered within its specified range. A fault-injection
technique deliberately violates one of these guarantees for a very short time, so
that a flip-flop latches a value that has not finished computing, or a memory
cell flips. Because the violation is brief and localised, most of the
computation stays correct and only a few bits are disturbed --- which is exactly
what differential fault analysis needs.

\section{Common injection techniques}
\begin{description}
\item[Clock glitching.] When the attacker controls the device clock, inserting
one abnormally short clock period forces a register to sample before the slowest
combinational path has settled. The bits on the slow paths latch incorrectly.
The resulting fault depends on the implementation and timing. In the
ChipWhisperer-Lite setup, clock injection requires an accessible target clock
path. No single-bit accuracy or completed physical STM32 campaign is established
by the software experiments in this report.
\item[Voltage glitching.] A brief, sharp drop (or spike) on the supply rail slows
gate switching so that, again, logic fails to settle in time. It needs no clock
access and little equipment --- a transistor pulling the rail low (a ``crowbar'')
is enough --- which makes it one of the most widely used techniques. Its spatial
and bit-level precision depends on the target and experimental arrangement;
neither method guarantees a selected bit flip.
\item[Electromagnetic (EM) pulse injection.] A small coil held over the package
emits a fast magnetic pulse that induces transient currents in nearby logic. It
is contactless, needs no decapsulation, and can be aimed at a region of the die,
giving some spatial selectivity.
\item[Optical / laser injection.] A focused light or laser pulse on an exposed
(decapsulated) die deposits charge in specific cells. It is the most precise
technique and can come closest to a true, chosen single-bit fault, but it is
invasive and requires expensive equipment and sample preparation.
\item[Auxiliary stressors.] Temperature and body-biasing shift timing margins and
are sometimes combined with the above to widen the window in which a glitch
succeeds.
\end{description}

\section{Hitting the right instruction}
Producing a disturbance is only half the problem; it must land on the targeted
operation (for Giraud, the byte just before the final SubBytes). Two controls do
this:
\begin{itemize}
\item A \textbf{trigger} marks when the sensitive computation runs. In this
project the firmware raises a GPIO line (TIO4) around the cipher, and the capture
hardware watches it. A real evaluation may instead trigger on a power or EM
feature.
\item A \textbf{delay} after the trigger selects \emph{which} clock cycle, and
therefore which instruction, receives the glitch; the glitch \emph{width} and
\emph{offset} set its strength. An evaluation sweeps these parameters, repeats
each setting many times, and classifies the outcome.
\end{itemize}

\section{The gap between the model and reality}
A real campaign does not produce clean single-bit faults on demand. Each attempt
typically falls into one of several outcomes:
\begin{center}
\begin{tabular}{ll}\toprule
Observed outcome & Use for DFA \\\midrule
No effect (correct output) & none \\
Reset or crash & none \\
Instruction skip / control-flow change & different attack class \\
Multi-bit or multi-byte fault & usually breaks the single-fault model \\
Clean one-bit / one-byte transient fault & \emph{this} is what Giraud/Piret need \\\bottomrule
\end{tabular}
\end{center}
\noindent The clean case is often a minority of attempts. The practical workflow
is therefore statistical: log the correct ciphertext and every faulty ciphertext
with its glitch parameters, \emph{filter} for the faults whose difference matches
the assumed model, and feed only those to the solver. This is why the software
experiments in this report separate the target (which may expose internal state)
from the solver (which sees only ciphertexts): the solver is written to accept
exactly the observations a real, filtered campaign would yield.

\note{The numerical examples elsewhere in this report inject an exact XOR in
software, so the fault model holds by construction. On the physical STM32 target
the same model must be \emph{established experimentally}: knowledge of the exact
bit (the \texttt{Giraud-known} variant) is a strong assumption that a campaign has
to justify by characterising its own glitches.}

\section{A taxonomy of fault widths and their consequences}
\label{sec:fault_taxonomy}
A disturbed computation can land anywhere on a spectrum, from no effect to a
corrupted whole word. What decides whether Giraud or Piret can use it is, above
all, the \emph{width} of the fault --- how many bits or bytes actually changed.
The two attacks are defined by the width they expect:
\begin{itemize}
\item \textbf{Giraud} expects \emph{exactly one bit} flipped at the input of the
last S-box layer.
\item \textbf{Piret} expects \emph{exactly one byte} corrupted (any nonzero
value) before the round-9 MixColumns.
\end{itemize}

\paragraph{Consequence by fault width.}
This is the dimension that decides whether an attack \emph{works} or \emph{breaks}.
\begin{center}
\begin{tabular}{P{.26\linewidth}P{.33\linewidth}P{.33\linewidth}}\toprule
Fault width & Giraud (model: 1 bit) & Piret (model: 1 byte) \\\midrule
Zero bits (no effect) & \(C=C'\), useless & \(C=C'\), useless \\
One bit & \emph{ideal} & works --- one bit is a byte of Hamming weight one \\
2--8 bits in the \textbf{same} byte & breaks: input difference has weight \(>1\), the weight-one filter rejects the true key & \emph{ideal} --- this is exactly a single-byte fault \\
One whole byte & breaks (same reason) & \emph{ideal} \\
\(\ge 2\) different state bytes & breaks the single-fault model & breaks: pattern no longer matches one byte in a column \\\bottomrule
\end{tabular}
\end{center}
\noindent The central observation is that \textbf{Piret is a superset of Giraud}:
every one-bit fault is also a one-byte fault (a byte value of Hamming weight one),
so the Piret solver accepts a bit, whereas Giraud demands weight one and rejects
anything wider. This is why, on SM4, the two attacks nearly coincide --- the only
difference is that Giraud assumes \(e\) is a single bit while Piret accepts the
whole byte.

\paragraph{Consequence by fault position (unknown location).}
Not knowing \emph{where} the fault landed does \emph{not} break the attack; it is
recovered by observation and raises the cost.
\begin{center}
\begin{tabular}{P{.3\linewidth}P{.36\linewidth}P{.26\linewidth}}\toprule
Situation & How it is resolved & Consequence \\\midrule
Unknown which column / key group & the set of changed ciphertext positions is the signature (see the column\(\to\)positions table of the Piret chapter) & none --- the group is \emph{read} from \(C\xor C'\) \\
Unknown exact position within the group & the solver sweeps all rows: \(\mathcal D=\{M(e\,\mathbf e_r):r\in\{0,1,2,3\},\,e\in\{1,\ldots,255\}\}\) & weaker filter \(\Rightarrow\) more faults \\\bottomrule
\end{tabular}
\end{center}
\noindent For AES Giraud this is even more direct: a one-bit fault changes exactly
one ciphertext byte (no MixColumns in the last round), so the changed position
already names the attacked key byte.

\paragraph{Consequence by fault value (unknown pattern).}
Not knowing the exact value does not break the attack either; it only weakens the
filter.
\begin{center}
\begin{tabular}{P{.16\linewidth}P{.27\linewidth}P{.27\linewidth}P{.14\linewidth}}\toprule
Attack & Value known & Value unknown & Measured \\\midrule
SM4 Giraud & exact bit: one mask & unknown bit: eight masks & median 32 versus 35 \\
AES Piret & \multicolumn{2}{P{.54\linewidth}}{The AES differential model can enumerate unknown nonzero byte errors.} & model dependent \\
SM4 controlled DFA & exact byte mask & not implemented by this known-mask filter & no claim \\\bottomrule
\end{tabular}
\end{center}
\noindent Distinguish the AES model from the project's SM4 controlled DFA:
the latter needs the exact error mask. The two plotted SM4 Piret configurations
change the ordered masks, not whether their values are known. Hiding an arbitrary
byte error would require a different filter and a new experimental evaluation.

\paragraph{Consequence by timing (wrong round or instruction).}
If the glitch lands on the wrong round or instruction, the difference propagates
through a different number of S-box layers and the pattern of changed bytes no
longer matches the expected one:
\begin{itemize}
\item Too early $\Rightarrow$ extra S-box layers scramble the difference $\Rightarrow$ matches no pattern $\Rightarrow$ filtered out.
\item Too late (e.g.\ already at the output) $\Rightarrow$ wrong positions change $\Rightarrow$ filtered out.
\end{itemize}
The delay sweep of \S\ref{sec:fault_taxonomy}'s sibling section exists precisely
to place the glitch on the right cycle.

\paragraph{Summary.}
\begin{center}
\begin{tabular}{P{.32\linewidth}P{.17\linewidth}P{.15\linewidth}P{.2\linewidth}}\toprule
What you do not know / what happened & Giraud & Piret & Action \\\midrule
No effect & useless & useless & discard \\
One bit, bit known & \checkmark\ 32 faults & \checkmark\ works & solve \\
One bit, bit unknown & \checkmark\ \({\approx}34.6\) & \checkmark\ works & sweep \(e\) \\
Several bits, same byte & breaks & \checkmark\ ideal & Piret solver \\
One byte, pos.\ unknown & n/a & \checkmark\ more faults & read group, sweep \(r\) \\
Several bytes & breaks & breaks & discard \\
Wrong round / instr. & no match & no match & retune delay \\\bottomrule
\end{tabular}
\end{center}
\noindent The rule that ties it together: the less you know (value, position), the
\emph{weaker the filter} and the \emph{more faults} you need --- but the attack
stays valid. What actually \emph{breaks} an attack is the width leaving the model:
more than one byte breaks both, while more than one bit breaks only Giraud (Piret
still holds).

\section{How can one bit actually be changed?}
In the simulator, $a'=a\xor2^j$ explicitly flips bit $j$. On the STM32F303,
the host generally cannot command ``flip this internal bit'' through the
ChipWhisperer-Lite. Its clock and voltage pulse generators perturb physical
conditions. A selected timing window may produce no change, a reset, an
instruction fault, or a data error of varying width. Laser and EM injection
require additional equipment; spatial targeting still does not by itself prove
that exactly one desired algorithmic bit was flipped.

For a diagnostic firmware build on the laboratory target, compare the same
intermediate word in correct and faulted executions: $e=A\xor A'$.
If $e\ne0$ and $e\mathbin{\&}(e-1)=0$, its Hamming weight is one.
For example $\hex{96FF9D45}\xor\hex{97FF9D45}=\hex{01000000}$.
Instrumented intermediate states are validation aids and must not be given to
the ciphertext-only solver. Instrumentation also changes timing, so a result
must be revalidated for the evaluated firmware build.

Without intermediate-state access, unchanged final SM4 words and a one-byte
$L^{-1}(\Delta X)$ are useful consistency checks. They do not prove that the
S-box input error had Hamming weight one. Candidate filtering tests compatibility
with that assumption; independent encryptions and fresh validation pairs should
check a recovered key. Knowledge of the exact mask requires evidence beyond
knowing the pulse parameters.

The hardware capabilities are documented by NewAE in
\href{https://chipwhisperer.readthedocs.io/en/latest/Capture/ChipWhisperer-Lite.html}{the ChipWhisperer-Lite reference}.
This establishes available pulse-generation mechanisms, not a guarantee of a
particular single-bit fault on the STM32F303.

\section{Countermeasures}
Because the attacks depend on obtaining a clean faulty ciphertext that matches a
known model, the defences aim to deny exactly that:
\begin{description}
\item[Redundancy and verification.] Compute the cipher twice, or decrypt the
result and compare; release the output only if both agree. A single transient
fault is then detected and suppressed.
\item[Infection.] If a difference is detected, scramble the output so the faulty
ciphertext carries no usable single-byte relationship instead of the clean
difference the solver needs.
\item[Sensors and detectors.] Clock, voltage, EM and light sensors reset the
device when an injection is attempted.
\item[Randomisation.] Masking and random timing make it harder to aim a fault at
a specific intermediate value and to interpret the resulting difference.
\end{description}
\noindent None of these is implemented in this project's deliberately
unprotected educational target; their absence is what makes the target suitable
for measuring the attacks in the first place.
